\documentclass[12pt]{article}
\usepackage[utf8]{inputenc}
\usepackage[T1]{fontenc}
\usepackage{hyperref}
\title{Arpitan Language Resources for K-12 Curriculum Development}
\author{Compiled by Claude Sonnet 5.5 (claude-sonnet-5-5) for Andr\'as Kornai}
\date{}
\begin{document}
\maketitle

\section{Introduction}
Arpitan (Franco-Proven\c{c}al) is spoken in France, Italy (notably the Aosta Valley) and Switzerland. Searches found general and orthography-related pages but no school textbooks or confirmed teaching programmes.

\section{Primary Educational Resources}
None found in search.

\section{Linguistic and Reference Materials}
The ORB common orthography by Dominique Stich (document on the ACA site): \url{https://www.arpitania.eu/aca/documents/Stich/ORB%204.9.18.docx}.

A Swiss site discussing the name Arpitan and ORB: \url{https://arpitan.ch/spip.php?article361}.

Omniglot overview of writing systems and orthographies: \url{https://omniglot.com/writing/francoprovencal.htm}.

Language profile (Universit\'e Laval): \url{https://axl.cefan.ulaval.ca/monde/franco-provencal.htm}.

\section{Religious and Community Resources}
Overview of Savoy and Aosta minority context: \url{https://www.eurominority.eu/index.php/savoy-and-aosta/}.

\section{Digital and Audio Resources}
A Google Summer of Code 2020 project page that discusses Arpitan standards (ORB chosen): \url{https://summerofcode.withgoogle.com/archive/2020/projects/6217740206997504}.

\section{Curriculum Development Context}
Standardisation is contested (several orthographies are in use), and no ready-made school materials were found in these searches. Further targeted searches in Aosta Valley, Savoy and Swiss sources would be needed.

\end{document}